\section{An explicit coefficient-three refinement with five retained labels}
\label{sec:coefficient-three-refined}

\paragraph{A refinement of a sufficient family.}
We keep the original rings, discrete coefficient and full strict condition
of (G177.1)--(G177.2).  A direct numerator calculation improves the
high-conductor threshold, and exact finite special values add several
lower-conductor labels.  The five labels left outside the resulting set
are not asserted to have nonzero strict groups.

Let $\mathscr D$ be the following disjoint set of nontrivial primitive
Dirichlet characters, subsequently restricted to $G_K$:
\[
 \begin{split}
 \mathscr D={}&\{\chi_{-4},\chi_8\}\\
 &\sqcup\{\delta:\operatorname{cond}(\delta)=16\}\\
 &\sqcup\{\delta:\operatorname{cond}(\delta)=32,\ \delta(-1)=-1\}\\
 &\sqcup\{\delta:\operatorname{cond}(\delta)=2^c,\ c\geq6\}.
 \end{split}\tag{G178.1}
\]

\begin{theorem}[AC-G178: refined original-coefficient strict zero at three]
\label{AC-G178}
For every $\delta\in\mathscr D$, every original finite coefficient
extension $E/\mathbb Q_3$ whose ring of integers $O$ contains its values,
every $e\geq1$, $s\in\mathbb Z$, and continuous
$\kappa:G_K\to1+\mathfrak m_O$ unramified away three, one has
\[
 \mathcal S_3\bigl(K,(E/O)
          (\delta\,\omega_3^{-s}\kappa^{-1}c_3^{-e})\bigr)=0.
 \tag{G178.2}
\]
The residual strict group is zero as well.  The coefficient ring can
have arbitrary absolute ramification.

For an actual page $h\geq2$, the number of nontrivial pure labels from
(G178.1) with conductor dividing $2^{h+2}$ is
\[
 \begin{aligned}
 D_3(h)&=6+4\,\mathbf1_{h\geq3}
                 +\max\{0,2^{h+1}-16\},\\
 D_3(h)&=2^{h+1}-6\qquad(h\geq3).
 \end{aligned}\tag{G178.3}
\]
Every such label has strict zero at every positive weight on the original
actual $O_3$.  Adding the previously proved trivial label gives
$2^{h+1}-5$ pure two-power labels when $h\geq3$.
The five labels not supplied by this union are precisely $\chi_{-8}$
and the four even primitive characters of conductor $32$.
This describes the proved set, not a maximal vanishing classification.
\end{theorem}

\begin{proof}
\proofstep{1. The exact restricted numerators.}
For $c\geq6$ write $\delta=\chi_{-4}^{\,a}\psi$, with $\psi$ even
and primitive of conductor $2^c$.  The only odd fixed factors
$\chi_{-4}^{\,a}\chi_{-3}^{\,b}$ are $\chi_{-4}$ and $\chi_{-3}$.
Work over the fixed field $k=\mathbb F_9$, with $n_0=3$.
Use the measure, factor $1/2$ special-value identity and one-zero
trace argument (G173.3)--(G173.14), not the stronger
minimal-polynomial-degree hypothesis of (G173.2).
Those steps require
\[
 c\geq2n_0=6,\qquad c>v_2(4f_\theta),\qquad
                         c>v_2(\operatorname{lcm}(f_\theta,8)).
 \tag{G178.4}
\]
All hold: $f_\theta$ is $4$ or $3$, respectively.
In particular every filter twist of the chosen primitive $2^c$th root
has full order and avoids the exceptional denominator roots.

If $L(0,\theta\psi)$ reduced to zero, the trace argument would give
$P_1(\zeta)=0$ for a primitive $2^c$th root $\zeta$, with precisely
the polynomial of (G173.14).  Its two possibilities are
\[
 \begin{array}{c|c|c|c}
 \theta&L& P_1(Z)&1-Z^L\\ \hline
 \chi_{-4}&8&Z&1-Z^8\\
 \chi_{-3}&24&Z-Z^{17}=Z(1-Z^{16})&1-Z^{24}.
 \end{array}\tag{G178.5}
\]
For the second row the indices are $1,9,17$; their character values
are $1,0,-1$.  For the first row the only index is $1$.
The denominators and numerators in (G178.5) are nonzero at
every primitive $2^c$th root for $c\geq6$.
This contradicts the forced zero.  Integrality was already established
by the exact Gauss-sum identity (G173.7), whose denominators are
three-adic units.  Thus all required odd products and their inverses
have three-adic unit values at zero for $c\geq6$.
This is a new explicit use of the previously proved trace identity;
it does not weaken a hypothesis of a cited theorem without proof.

\proofstep{2. Exact small-conductor values.}
For an odd primitive character $\nu$ of conductor $f$, calculate
\[
                  L(0,\nu)=-\frac1f\sum_{a=1}^{f}a\nu(a).
 \tag{G178.6}
\]
For $\delta$ of conductor $2^c$, $c=4$ or $5$, write
$\beta=\delta(5)$, a primitive $2^{c-2}$th root.  Each odd integer
has a unique expression $(-1)^j5^t$ modulo $2^c$.
Substitution in (G178.6), with $\beta^2=-1$ at $c=4$ and
$\beta^4=-1$ at $c=5$, gives the following exact identities:
\[
 \begin{array}{c|c|c|c}
 \delta&\nu&\sum_{a=1}^{f}a\nu(a)&L(0,\nu)\\ \hline
 \chi_{-4}&\chi_{-4}&-2&1/2\\
 \chi_8&\chi_8\chi_{-3}&-48&2\\
 c=4,\ \delta\text{ even}&\delta\chi_{-3}&96\beta&-2\beta\\
 c=4,\ \delta\text{ odd}&\delta&-16(1+\beta)&1+\beta\\
 c=5,\ \delta\text{ odd}&\delta&
       -32(1+\beta)(1-\beta^2)&(1+\beta)(1-\beta^2).
 \end{array}\tag{G178.7}
\]
Here the conductors $f$ of the five rows are $4,24,48,16,32$.
The displayed sums can also be obtained by summing
$a(-1)^{j}\beta^t$ for odd $\delta$, or
$a\beta^t\chi_{-3}(a)$ for even $\delta$, over the listed conductor.
All entries are identities in the indicated cyclotomic integer ring,
before any reduction.
In the two rows with $3\mid f$, the numerator is divisible by three
in that ring.  The factor three is cancelled integrally before reduction;
no division by zero in a residue field is used.

Every value in the last column is a three-adic unit at every
coefficient embedding.  Indeed roots of order four or eight retain
their exact order in characteristic three, so $\beta$, $1+\beta$ and
$1-\beta^2$ appearing there are nonzero.  The same argument holds
with $\beta^{-1}$ in place of $\beta$.
Thus both inverse special values needed for the two rational
extensions of each allowed low-conductor label are units.
The exact polynomial sums and independent residue-field evaluations
are retained as finite checks; the factorization just given proves
the assertion at all embeddings.

\proofstep{3. Primitive integral annihilation and reflection.}
For an allowed character of conductor $2^c$, put
$H=\mathbb Q(\zeta_{3\cdot2^c})$, $J=\operatorname{Gal}(H/K)$ and
$A=\operatorname{Cl}(H)\{3\}$, exactly as in (G177.8).
Here $c\geq2$, $[H:\mathbb Q]=2^c$ and $[H:K]=2^{c-1}$.
Both degrees are prime to three.
The rational extensions of the residual character are $\delta$
and $\delta\chi_{-3}$.
For an odd extension $\nu$, use its primitive conductor $f_\nu$
and the integral Stickelberger annihilator (G177.9).
For all allowed $c\geq3$ choose
$b_0\equiv5\pmod{2^c}$ and $b_0\equiv1\pmod3$.
Then $\delta(5)$ has order $2^{c-2}>1$.
For $\delta=\chi_{-4}$ instead use $b_0=7$.
In every case $b_0$ is prime to $f_\nu$, reduces to one modulo three,
and $\nu(b_0)=\delta(b_0)\ne1$ after reduction.
The same statements hold for inverse characters.
Consequently
\[
             -(b_0-\nu(b_0))L(0,\nu^{-1})
                         \quad\hbox{is a three-adic unit}.
 \tag{G178.8}
\]
For the two quadratic rows one can take $b_0=7,13$ respectively;
their scalars are $-4$ and $-28$, both units.
The integral annihilator therefore kills the corresponding odd
three-class component in the primitive-conductor field.
Norm and ideal extension to $H$, whose intervening degree is a
power of two, transfer this zero component exactly as in G177,
without an imprimitive Euler-factor assumption.

For completeness, the even component uses the injection from
(G177.11): its mirror has character $\omega_3\nu^{-1}$ and is odd.
Every completion of $H$ at three has ramification index two over
$\mathbb Q_3$, since the two-power cyclotomic extension is unramified
there.  A nonzero $\mu_3$ radical would adjoin $\mu_9$ with relative
ramification index three, so it cannot be an unramified radical.
The remaining odd unit part is zero in a CM field after removing
torsion.  The odd kernel of that sequence is therefore zero for these
low conductors as well.
The odd mirror components are precisely among the unit-value
components of Steps 1--2, including their inverses.
It follows that the $\delta$ and $\delta\chi_{-3}$ components of
$A$ are zero, regardless of their parity.
This is the same previously audited Stickelberger and reflection
argument of G177 Steps 2--3, with its hypotheses now checked for
every added conductor.

\proofstep{4. Residual strictness and the original local lift.}
On $G_K$ the residual coefficient is $\overline\delta$ for every
$e,s,\kappa$, because $\omega_3|_{G_K}=1$.
Restriction over the prime-to-three extension $H/K$, followed by
class field theory, injects the residual strict group into
\[
  \operatorname{Hom}_{\mathbb F_3[J]}
                  (A/3A,k'(\overline\delta)),
 \tag{G178.9}
\]
as in (G177.12)--(G177.13).
Both rational class components have just been killed.
Faithful finite residue-field extension, if needed to split the
characters, proves that (G178.9) is zero over the original residue field.

At the full local group $G_{K_3}$ the residue degree is one.
For every allowed $c\geq4$, the value $\delta(3)$ has exact order
$2^{c-2}>2$ by the $-3$ generator argument (G177.14).
For the two additional quadratic labels the values are
\[
                    \chi_{-4}(3)=\chi_8(3)=-1.
 \tag{G178.10}
\]
In every case the original scalar at a local Frobenius lift minus
one is an $O$-unit.  Therefore the local and global invariants
of $D$ vanish.
Away three the remaining finite inertia has two-power order.
Its invariants are either $D$ or zero, and hence divisible by the
original uniformizer $\varpi$.
Apply the exact sequences for
$0\to D[\varpi]\to D\xrightarrow{\varpi}D\to0$
globally and at all the stated local groups.
Their local maps from residual degree-one cohomology are injective
by this invariant calculation.  Thus any strict class killed by
$\varpi$ lifts to a residual strict class, which is zero.
Every class is killed by a power of $\varpi$, since a continuous
cocycle into the discrete torsion module has finite image.
This proves (G178.2) without extending or replacing the original lattice.

\proofstep{5. Actual labels, counts and precise exceptions.}
The same host and coefficient argument as G177 Step 6 applies:
all these characters of conductor dividing $2^{h+2}$ occur in
the actual host, their orders divide $2^h$, and restriction to
$G_K$ does not identify distinct two-power characters.
Inflation preserves the original full local and inertia conditions.
There are two named quadratic labels, four primitive conductor-16
characters, and four odd primitive conductor-32 characters.
The number of primitive characters at each $c\geq6$ is $2^{c-2}$.
Their geometric sum gives exactly (G178.3).
The previous G177 family has $C_3(h)=\max(0,2^{h+1}-64)$ labels.
For $h\geq6$ the present theorem therefore adds exactly
\[
                      D_3(h)-C_3(h)=58.
 \tag{G178.11}
\]
All these statements hold at every positive weight.

The list of excluded labels follows by counting all
$2^{h+1}$ characters of $(\mathbb Z/2^{h+2}\mathbb Z)^\times$,
including the trivial one, when $h\geq3$.
The old trivial-label assertion is AC-K11, with the original-$O_3$
transport (EW.8)--(EW.9) of AC-G146.
For $\chi_{-8}$ the local value at three is one, so the preceding
local injection argument does not apply.
For an even primitive conductor-32 character, (G178.6) instead gives
\[
 \begin{aligned}
 L(0,\delta\chi_{-3})&=2+2\beta^2-2\beta^3,\\
 \overline{L(0,\delta\chi_{-3})}
     &=(\beta+1)(\beta^2+\beta-1),\\
 \beta^4+1&=(\beta^2+\beta-1)(\beta^2-\beta-1)
                          \quad\text{in characteristic three}.
 \end{aligned}\tag{G178.12}
\]
Exactly one of the two quadratic factors vanishes at a primitive
eighth root; inversion exchanges them.
Thus at every coefficient embedding one of the two needed values,
for $\delta$ and $\delta^{-1}$, fails to be a unit.
There is no justified half-embedding addition to the strict family.
A nonunit annihilator, or failure of this invariant argument,
does not prove that the associated strict group is nonzero.
\end{proof}

\paragraph{Source and geometric scope.}
The primary trace inputs are Sinnott, Lemma 3.6, printed pp.215--217,
and Proposition 4.6, printed pp.220--222, in the complete specialization
already proved in G173.
The integral annihilator and reflection sources remain Sharifi,
Definition 4.5.1, Lemma 4.5.2, Theorem 4.5.6, Lemma 1.4.6,
Proposition 1.3.2 and Theorem 4.4.2 in the exact G177 source scopes.
The finite-cover and Kummer dictionary is the already audited
SGA 1, Expos\'e V, sections 4 and 7, and Stacks 03PK.
No conductor-seven, auxiliary quadratic or nontrivial-cubic factor
is supplied here.  The actual numerical remainder, mixed realization
$\mathcal R_{\rm cub}$, H04.4 and cubic RH remain open.
